\section{The recognition source}
\label{sec:recognition_source}

By Theorem~T, the Riemann hypothesis is the statement \(\AZ=0\).  The
duality-confinement route asks for this vanishing to come from
domination records of vanishing trace.  The following hypothesis
states what the conditional theorem of \Cref{sec:landing_chain}
consumes; the formal record carries in addition some structure that no
proof uses (weights, a response involution, a quantitative separation
modulus), listed in \Cref{app:formalization}.  It is phrased with an auxiliary
ledger in the abstract form of the duality-confinement theorem
\citep[Remark~7.2]{TsiokosSDTC2026}, because the formal statement
allows the domination records to live in any ordered set with a trace,
not only among operators.

\begin{obligation}[The recognition source \(\GamSDTCSelberg\)]
\label{obl:rh:gamma-sdtc-selberg}
Let \(\Sel\) be a shell.  We say that the recognition source
\(\GamSDTCSelberg\) \emph{holds for} \(\Sel\) if there exist the
following data.
\begin{enumerate}[label=\textup{(\roman*)}]
  \item \emph{An auxiliary ledger.}  A finite set \(X\) with an involution \(J\) (the map
    \(\rho\mapsto x_\rho\) of item (iv) need not be injective), a set of \emph{visible} objects, a response set
    \(Y\) with a zero element, and a readout
    \(\psim:X\to Y\) that is separating on visible objects:
    for visible \(x\), \(\psim(x)=0\) if and only if \(Jx=x\).
  \item \emph{An ordered set with a trace.}  A set \(C\) with a
    relation \(\preceq\), a subset of positive elements, a zero element,
    a distinguished positive element \(\AX\), and a map
    \(\tau:C\to\mathbb R\) that is monotone for \(\preceq\) and
    nonnegative on positive elements, with \(\tau(\AX)=0\) only if
    \(\AX\) equals the zero element.
  \item[(ii\('\))] \emph{A link between the two.}  If \(\AX\) is the zero
    element, then \(\psim(x)=0\) at every visible \(x\).
  \item \emph{Domination records of vanishing trace.}  Positive
    elements \(B_n\in C\) with \(\AX\preceq B_n\) for all \(n\) and
    \(\tau(B_n)\to0\).
  \item \emph{A bridge to the zeros.}  For each represented zero
    \(\rho\) of \(\Sel\), a visible object \(x_\rho\in X\), and a map
    \(e:R\to Y\) on the shell's coordinates with \(e(r)=0\) only if
    \(r=0\), such that
    \[
      \psim(x_\rho)=e\Bigl(c(\rho)-\frac12\Bigr).
    \]
\end{enumerate}
\end{obligation}

Items (i)--(iii) are an instance of the duality-confinement setting in
its abstract form: the domination records and the trace squeeze force
\(\AX\) to vanish, and item (ii\('\)) then gives \(\psim=0\) at every visible
object.  None of this involves the zeta function, and neither \(\AZ\) nor \(q\)
appears in the hypothesis; the zeros enter only through item (iv), which asserts that the auxiliary readout at the
object attached to \(\rho\) encodes the displacement coordinate
\(c(\rho)-\frac12\).  In this abstract form the link between
\(\AX\) and \(\psim\) is part of the data (item (ii\('\))) rather than a
consequence of a construction, and the next proposition shows that the
hypothesis is exactly as strong as the critical-line statement for the
represented zeros.  The operator form, in which \(\AX\) is the Gram
operator of the readout, is \Cref{thm:rh:concrete-window-dc}.  Nothing in the
hypothesis requires the auxiliary ledger to be constructed from
zeta-function data, and its intended provenance---an instance of the
companion paper's structural law obtained from a trace formula---is
not established here.

The next proposition measures the strength of the hypothesis.

\begin{proposition}[Calibration of the recognition source]
\label{prop:rh:gamma-calibration}
For every shell \(\Sel\), \(\GamSDTCSelberg\) holds for \(\Sel\) if and
only if \(c(\rho)=\frac12\) for every represented zero
\(\rho\).  In particular, \(\GamSDTCSelberg\) holds for the bare shell if and only
if the Riemann hypothesis holds.
\end{proposition}

\begin{proof}
If \(\GamSDTCSelberg\) holds, \Cref{thm:rh:conditional} below gives
\(c(\rho)=\frac12\) for every represented zero.
Conversely, suppose \(c(\rho)=\frac12\) for all represented
zeros.  Take \(X\) a single visible point with \(J\) the identity,
\(Y=\mathbb R\), \(\psim=0\), \(C=\mathbb R\) with the usual order and positive elements
\([0,\infty)\),
\(\AX=0\), \(B_n=0\), \(\tau\) the identity, \(x_\rho\) the point, and \(e:R\to\mathbb R\) with \(e(0)=0\) and \(e(r)=1\) for \(r\neq0\)
(a classical case distinction on the abstract set \(R\)).  Items (i)--(iii) hold trivially, \(e(r)=0\) only for \(r=0\), and item (iv)
holds because \(c(\rho)-\frac12=0\) by the coordinate
laws, so \(e(c(\rho)-\frac12)=0=\psim(x_\rho)\).  On the
bare shell, \(c(\rho)=\Real\rho\) and the represented zeros
are the nontrivial zeros, so the condition is the Riemann hypothesis
(\Cref{thm:rh:translation-T}).
\end{proof}

Thus the recognition source, as a type of data, is exactly as strong as
the critical-line statement it is used to prove, and on the bare shell
establishing it is no easier than proving the Riemann
hypothesis.  Its intended content is not its logical
strength but its provenance: the conditional theorem becomes a proof
only if the data are produced by an independent argument, for instance
from a trace formula, in which item (iv) is a consequence rather than
an assumption.  \Cref{thm:rh:concrete-window-dc} shows that the same
calibration holds for a non-degenerate version built from the actual
zeros.
